\section{Questions}\label{sec:questions}

\begin{enumerate}
\item Is every block map commuting with $F_{95}$ a constant or of the form $\widehat{H_{\varphi,\psi}}\sh^n$ with block maps $\varphi,\psi$ of the golden mean shift that conserve the number of particles (Remark~\ref{cl:rem:95end}(1))? This holds for the block maps of radius at most $4$, and for those of sublattice type of radius $5$ (Proposition~\ref{cl:prop:95census}).
\item Describe the group $\mathcal N$ of Section~\ref{cl:sec:95}: generators, abelianisation, centre; is it finitely generated, and does it contain a free group? By Lemma~\ref{cl:lem:95flux} it is the stabiliser in $\Aut(\mathsf G,\sh)$ of the class of the function $x\mapsto x_0$ modulo continuous coboundaries. (The boxes of Lemma~\ref{cl:lem:95box} carry one particle each, which suggests that the marker constructions of \cite{BLR1988} for free products of finite groups can be adapted to $\mathcal N$; this is a suggestion and not a claim.)
\item Determine $\Aut(X_{F_{91}})$. Lemma~L holds at every radius; what is missing is a normalisation of the automorphisms as in Theorem~\ref{cl:thm:KR94}, which has to use points that are not $3$-sparse (Proposition~\ref{cl:prop:91KR} and Remark~\ref{cl:rem:91}), and a proof that a normalised automorphism fixes every point of $\Phi^+\cup\Phi^-$.
\item Does Lemma~L hold for $155$ at every radius? It would follow from the front hypothesis (Fr) and the closure (Cl) of Remark~\ref{cl:rem:155G1}.
\item Is the class of $95$ the only class of elementary rules that are not affine for which $F_f$ has no blocking word and $\Aut(X_{F_f})$ is not virtually $\Z^2$? This is Question~1.6 of \cite{PaperC}; by the classification of \cite{PaperD} it concerns the classes whose automorphism group is not known. For the class of $175$, whose rules also do not read their middle argument, see Remark~\ref{cl:rem:175}.
\end{enumerate}
